%  LaTeX support: latex@mdpi.com 
%  For support, please attach all files needed for compiling as well as the log file, and specify your operating system, LaTeX version, and LaTeX editor.

%=================================================================
\documentclass[electronics,article,submit,pdftex,moreauthors]{Definitions/mdpi} 
% For posting an early version of this manuscript as a preprint, you may use "preprints" as the journal and change "submit" to "accept". The document class line would be, e.g., \documentclass[preprints,article,accept,moreauthors,pdftex]{mdpi}. This is especially recommended for submission to arXiv, where line numbers should be removed before posting. For preprints.org, the editorial staff will make this change immediately prior to posting.

\graphicspath{{figures/}} % path to folder with figures
\usepackage{caption}
\usepackage{subcaption}
\usepackage{xcolor}
\usepackage{graphicx}
\usepackage{gensymb} % for degree symbol
\usepackage[super]{nth}
\usepackage{color, colortbl}
\definecolor{Goldenrod}{RGB}{255,228,181}
\usepackage{tabularx}
\usepackage{chemmacros}

%=================================================================
% MDPI internal commands
\firstpage{1} 
\makeatletter 
\setcounter{page}{\@firstpage} 
\makeatother
\pubvolume{1}
\issuenum{1}
\articlenumber{0}
\pubyear{2022}
\copyrightyear{2022}
%\externaleditor{Academic Editor: Firstname Lastname}
\datereceived{} 
%\daterevised{} % Only for the journal Acoustics
\dateaccepted{} 
\datepublished{} 
%\datecorrected{} % Corrected papers include a "Corrected: XXX" date in the original paper.
%\dateretracted{} % Corrected papers include a "Retracted: XXX" date in the original paper.
\hreflink{https://doi.org/} % If needed use \linebreak
%\doinum{}
%------------------------------------------------------------------
% The following line should be uncommented if the LaTeX file is uploaded to arXiv.org
%\pdfoutput=1

%=================================================================
% Add packages and commands here. The following packages are loaded in our class file: fontenc, inputenc, calc, indentfirst, fancyhdr, graphicx, epstopdf, lastpage, ifthen, lineno, float, amsmath, setspace, enumitem, mathpazo, booktabs, titlesec, etoolbox, tabto, xcolor, soul, multirow, microtype, tikz, totcount, changepage, attrib, upgreek, cleveref, amsthm, hyphenat, natbib, hyperref, footmisc, url, geometry, newfloat, caption

%=================================================================
%% Please use the following mathematics environments: Theorem, Lemma, Corollary, Proposition, Characterization, Property, Problem, Example, ExamplesandDefinitions, Hypothesis, Remark, Definition, Notation, Assumption
%% For proofs, please use the proof environment (the amsthm package is loaded by the MDPI class).

%=================================================================
% Full title of the paper (Capitalized)
\Title{Simplex Lattice Design and X-Ray Diffraction for Analysis of Soil Structure: A Case of Cement-Stabilized Compacted Tills Reinforced with Steel Slag and Slaked Lime}

% MDPI internal command: Title for citation in the left column
\TitleCitation{Simplex Lattice Design and X-Ray Diffraction for Analysis of Soil Structure: A Case of Cement-Stabilized Compacted Tills Reinforced with Steel Slag and Slaked Lime}

% Author Orchid ID: enter ID or remove command
\newcommand{\orcidauthorA}{0000-0002-0577-9936} % Add \orcidA{} behind the author's name
\newcommand{\orcidauthorB}{0000-0002-5759-1089} % Add \orcidA{} behind the author's name

% Authors, for the paper (add full first names)
\Author{Per Lindh $^{1,2}$\orcidA{} and Polina Lemenkova $^{3}\orcidB{}$*}

%\longauthorlist{yes}

% MDPI internal command: Authors, for metadata in PDF
\AuthorNames{Per Lindh, Polina Lemenkova}

% MDPI internal command: Authors, for citation in the left column
\AuthorCitation{Lindh, P.; Lemenkova, P.}
% If this is a Chicago style journal: Lastname, Firstname, Firstname Lastname, and Firstname Lastname.

% Affiliations / Addresses (Add [1] after \address if there is only one affiliation.)
\address{%
$^{1}$ \quad Swedish Transport Administration, Malm\"{o}, Sweden; per.a.lindh@gmail.com \\
$^{2}$ \quad Lund University, Lunds Tekniska H\"{o}gskola LTH (Faculty of Engineering), Department of Building and Environmental Technology, Division of Building Materials, Box 118, SE-221-00, Lund, Sweden.\\
$^{3}$ \quad Universit\'{e} Libre de Bruxelles, \'{E}cole polytechnique de Bruxelles (Brussels Faculty of Engineering), Laboratory of Image Synthesis and Analysis, Building L, Campus de Solbosch, ULB -- LISA CP165/57, Avenue Franklin D. Roosevelt 50, B-1050, Brussels, Belgium}

% Contact information of the corresponding author
\corres{Correspondence: polina.lemenkova@ulb.be; Tel.: +32471860459 (P. Lemenkova)}

% Current address and/or shared authorship
% \secondnote{These authors contributed equally to this work.}
% The commands \thirdnote{} till \eighthnote{} are available for further notes

%\simplesumm{} % Simple summary

%\conference{} % An extended version of a conference paper

% Abstract (Do not insert blank lines, i.e. \\) 
\abstract{Evaluating the structure of soil prior to building construction has become more valuable in a large variety of geotechnical and civil engineering applications. To built an effective framework for assessment and modelling stabilized soil specimens, the workflow can include a complex approach of simplex lattice design and X-Ray diffraction for analysis of soil structure. Different from traditional in-situ measurements, we propose a statistical framework for effective decision on binder combination in the process of soil stabilization. Practical aim of this study was to evaluate strength properties of different soil types using ordinary Portland cement, slaked lime and steel slag as pure agents and blended binders. Soil specimens were collected in southern Sweden and include sandy silty tills and clay till (clay content 6-18\%). The pre-processing included mineralogical analysis of mineral composition and soil structure by X-ray diffraction (XRD) and sieve. The specimens were fabricated, compacted, rammed, stabilized by six binder blends and assessed for Uniaxial Compressive Strength (UCS). Moisture Condition Value (MCV) and water content tests were done for compacted soil and shown variation in the MCV values for different binders. The study determined the effects from binder blends on the UCS gain in three types of soil, measured on days 7, 28 and 90th. Positive effects were noted from steel slag/lime blends on UCS gain in sandy silty tills. Steel slag/slaked lime mixed binder performs better compared to pure binders. The effectiveness of simplex lattice design is demonstrated in a series fo ternary diagrams showing soil strength evaluated by adding stabilising agents in different proportions.}

% Keywords
\keyword{steel slag; simplex lattice design; ordinary Portland cement; slaked lime; X-Ray diffraction; compressive strength; soil stabilization; civil engineering} 

% The fields PACS, MSC, and JEL may be left empty or commented out if not applicable
%\PACS{J0101}
%\MSC{}
%\JEL{}

%%%%%%%%%%%%%%%%%%%%%%%%%%%%%%%%%%%%%%%%%%
% Only for the journal Diversity
%\LSID{\url{http://}}

%%%%%%%%%%%%%%%%%%%%%%%%%%%%%%%%%%%%%%%%%%
% Only for the journal Applied Sciences
%\featuredapplication{Authors are encouraged to provide a concise description of the specific application or a potential application of the work. This section is not mandatory.}
%%%%%%%%%%%%%%%%%%%%%%%%%%%%%%%%%%%%%%%%%%

%%%%%%%%%%%%%%%%%%%%%%%%%%%%%%%%%%%%%%%%%%
% Only for the journal Data
%\dataset{DOI number or link to the deposited data set if the data set is published separately. If the data set shall be published as a supplement to this paper, this field will be filled by the journal editors. In this case, please submit the data set as a supplement.}
%\datasetlicense{License under which the data set is made available (CC0, CC-BY, CC-BY-SA, CC-BY-NC, etc.)}

%%%%%%%%%%%%%%%%%%%%%%%%%%%%%%%%%%%%%%%%%%
% Only for the journal Toxins
%\keycontribution{The breakthroughs or highlights of the manuscript. Authors can write one or two sentences to describe the most important part of the paper.}

%%%%%%%%%%%%%%%%%%%%%%%%%%%%%%%%%%%%%%%%%%
% Only for the journal Encyclopedia
%\encyclopediadef{For entry manuscripts only: please provide a brief overview of the entry title instead of an abstract.}

%%%%%%%%%%%%%%%%%%%%%%%%%%%%%%%%%%%%%%%%%%
\begin{document}

%%%%%%%%%%%%%%%%%%%%%%%%%%%%%%%%%%%%%%%%%%
\section{Introduction}

Cementitious materials have positive effects on developing of strength and hardening of clayey soil during stabilization. Cement is the most widely used traditional binder often used for soil stabilization \cite{Moh,Lindh2021a,Lapointe,ZhangTao}. Apart from cement, many other stabilizing agents are being added to binder blends. Replacing the cement by other materials aims to enhance technical performance of stabilization, increase the economic effectivity of the construction process through cheaper materials and reduce environmental impacts from $CO_2$ discharges \cite{Joulazadeh,Yietal2014,Renjith}. Besides, the effects from blends fabricated from various binders on soil are more durable and notable \cite{Lindh2022b,Mariri,Fasihnikoutalab}. The stabilization of soil is owing to the ability of binders to decrease the space between the pores in soil through bonding of particles during the period of curing \cite{Lindh2021b,Bahrekazemi}. The decrease of volume in the pores of soil develops over the time of curing and results in the decreased porosity and permeability of soil \cite{Lindh2022b,Mariri,Fasihnikoutalab}. Thus, cementitious binders increase the compactness and strength of soil, which finally increases its workability and bearing capacity, – the key requirements for soil quality in road construction.	

Many various binders make notable beneficial effects on the development of strength in binder-soil blends in the process of stabilization. For instance, adding slaked lime and steel slag increases Uniaxial Compressive Strength (UCS) of soil, as reported in many studies \cite{Lindh2021b,Bahrekazemi,Lindh2022c}. As for the impact from the mixtures of binders on soil compactness, earlier works shown their positive effects on various types of soil intended for construction of roads \cite{Dahlin,Nordmark,Kallen}. Ground Granulated Blast Furnace Slag (GGBS) is a steel waste with pozzolanic properties, received as an industrial by-product from the steelmaking process \cite{Aziz}. 

The fundamental physical properties of metallurgical slag include viscosity and rheological parameters, formed through chemical composition which includes SiO2 (silica), Al2O3 (aluminium oxide) and MgO (magnesium oxide) \cite{Slezak}. This makes steel slag a widely used as binding component in soil stabilization for roadbed structures \cite{Mraz,Lindh2004,Erdenebold}, as it has beneficios impacts on pore structure of soil, due to its physico-chemical properties. Because soil as a porous construction material experience shrinkage and decrease in voids with the inclusion of cementitious binders, the compaction and stabilization processes significantly decrease porosity and increase bulk density of soil \cite{Bhattacharjee}. This suggests positive effect from additives in soil, which increase its strength and decrease permeability and porosity when added in optimal amounts \cite{Ozelim,Lindh}. Another often-studied environmental parameter that is closely related to climate setting is the freeze-thaw durability of soil used as roadbeds, which is actual for countries with cold climate conditions \cite{Su,Sahlabadi}. The UCS of soil remarkably increases with the inclusion of binders, and suggests beneficial effects from reinforcing materials on the resistivity of soil to the seasonal climate cycles, such as freeze-thaw processes. 

The chemistry of stabilizing agents should also be considered when dealing with the effects of binders on soil strength \cite{Bowers}. Earlier works reported that the microstructure and geotechnical parameters of soil (e.g., shear modulus, compaction, Atterberg limits, internal friction angle, and maximum dry unit weight) improve along with added hardeners, such as cement pastes or other stabilizing agents, mainly due to the increased bonds between soil particles, which chemically reinforces its inner structure \cite{Ghorbani,Rajabi,Zhangetal2018}. Steel slag, in particular, binds soil particles effectively in contaminated soil polluted by heavy metals, and decreases leaching \cite{Negim,Lindh2022d,Moon}. In its turn, steel mill sludge can be used for reinforcement of clay bricks, to decrease negative environmental consequences from its discharge and to optimise its utilization \cite{Sarani}.

The effects from blended binders and their influence on the development of soil strength are not as straightforward. While in general the inclusion of stabilizing agents increases the UCS of soil, the effects may vary individually depending on the types of binders and soil and external factors (curing period and temperature). For example, slaked lime and fly ashes reduce the permeability and improves the strength in the cohesive soil due to pozzolanic reaction from siliceous and aluminous material \cite{Sharma2012}. The incorporation of cement decreases the porosity in pore microstructure of soil due to the bonding of particles \cite{Inyang,Consoli}. However, the effects from the soil types and different pH values also increase strength capacity due to higher pH values \cite{Gori,Ghobadi}, as well as with regard to different soil types (fine grained, medium grained or coarse-grained) \cite{Malikzada}. Although rates in strength development do generally increase with increasing content of binders in cement-soil mixture, external factors (moisture content and temperature) also play a significant role. Thus, since the effects of pore refinement in soil and compaction act simultaneously during curing period, it is alway necessary to determine individually in each case the effects from binders that reduce permeability of soil, gain in strength and environmental effects (decreased leaching, reduced pH), which might vary along with the amount and the type of binders used for soil stabilization.

The existing literature and case studies on soil stabilization tended to focus on the final stage assessing the gain in strength. It can be the effects from various binders, both traditional \citep{Boswell,Sharma2017,Sahuetal2017} or novel binders \cite{Saffari,YadavTiwari}, and the evaluation of workability of stabilized soil \cite{Ogila}. These methods are effective to evaluate the general stabilization results and to obtain a rough estimate of the soil strength over a curing period. However, a comparative analysis of blended mixtures needs special attention, to provide adequate data regarding the analysis of the effects from particular binders and their ratios, when mixed in proportions, on soil strength. Apart from the traditional UCS testing used to evaluate the strength of stabilized soil, more advanced techniques that allow for accurate and non-destructive measurements of specimens include seismic waves and X-rays \cite{Kierlik,Wani}. One such technique is the propagating the elastic P-waves through the soil mass, where strength is calculated from the increase in wave velocities which directly correlate with the stiffness, hardness and density in microstructure of the porous media \cite{Fakhrabadi,Sharma2011}.

Accurate and systematic evaluation of strength in soil stabilized by various binders using blended mixtures in various proportions are lacking, especially those focused on Swedish clayey tills, which are subject to harsh environmental and climate conditions. The objective of this study is a comparative analysis of the effects from various binders and their proportions on soil strength over a period of curing. Besides, the research on influence of binders on Moisture Condition Value (MCV) of soil need to be performed to show their effects with aim to soil improvement. The non-selective use of binders for any type of soil can result in the opposite effects, since as various types of binders are better adjusted to the various soil types, respectively. Therefore, the effects from diverse binders should be tested and evaluated separately using statistical approach for optimization of the workflow. To this end, we performed a series of statistical simplex tests to define the best relationships of soil-binder mixtures with regard to the strength of soil.

The goal of this study was to determine the effects that stabilizing agents (ordinary Portland cement, slaked lime and slag GGBFS have on the long-term strength gain in 3 types of soil measured on days 7, 28 and 90th. Besides, other geotechnical characteristics were tested for each type of soil stabilized by various binders, respectively: dry density, water content and MCV. In order to perform tests more effectively, we performed statistical optimization of these processes using existing experimental approaches \cite{Montgomery,MyersMontgomery}. The research has been carried out in the Swedish Geotechnical Institute (SGI) using the available tools and equipment. The methodology followed the existing standards on soil quality testing defined in technical guidances by the Swedish Institute for Standards (SIS). The dynamics in strength and changes of soil structure stabilized by various binders were examined by the UCS tests, moisture and compactness were monitored with Proctor compaction technique during curing period. The resulting changes were visualised, compared and assessed in series of statistical ternary plots, presented below in the article. The correlations between binders and strength in the stabilized soil mixture, the estimated UCS response surfaces for different recipes of soil-binder mixtures and the observed values were detected, plotted and evaluated.
	
%%%%%%%%%%%%%%%%%%%%%%%%%%%%%%%%%%%%%%%%%%
\section{Materials and Methods}

\subsection{Approach}

To evaluate the structural and physicochemical properties of the soil, we applied two types of tests: the Moisture Condition Value (MCV) for compaction characteristics of specimens; and the Uniaxial Compressive Strength (UCS). These two parameters were analysed to evaluated the compaction properties and gained strength of the stabilized soil.

\subsection{Material characteristics}
The soil specimens used in this study were fabricated from the raw material originated from three different trial pits in southern Sweden. The soil included different types A, B and C with varied physio-chemical characteristics, described as follows. The clay content of soil types A, B and C varied from 6 to 18\%. Based on the dominating size of the particles within tested specimens, soils of types A and C are sandy silty tills, while soil type B is clay till. Soil type A had the following characteristics: particle density $\rho _s =2.70 mg/m^3$ ; Optimum Moisture Content (OMC) according to the modified Proctor is 6.1\%; dry density $\rho _d = 2.26 Mg/m^3$; natural water content (WN) =9.8\%; Liquid limit measured by cone penetrometer (WL) is 19.3\%, plasticity index <10. Soil type B had the following characteristics: particle density $\rho _s = 2.71 mg/m^3$; OMC as modified Proctor – 9.2\%; dry density $\rho _d = 2.12 Mg/m^3$; natural water content (WN) =14.6\%; Liquid limit by cone penetrometer (WL) 23.9\%, plasticity index <11.9. Soil type C, included only measurements of the particle density and WN. Thus, particle density $\rho _s$ for soil type C was $2.67 mg/m^3$ and WN =13.0\%; plasticity index <10. The summary of the performed tests on soil types A, B and C is presented in Table \ref{tab:1}. 

\begin{table}[H]
\caption{Summary of the test on soil types A, B and C* stabilised by cement (OPC), lime and slag (GGBFS)}
\small
    \centering
    \begin{tabular}{|l|l|l|l|l|l|l|l|l|l|l|l|}
    \hline
    	 \rowcolor{Goldenrod}
        Soil type A & ~ & ~ & ~ & ~ & ~ & ~ & ~ & ~ & ~ & ~ & ~ \\ \hline
        Recipe & 0 & 1 & 2 & 3 & 4 & 5 & 6 & 7 & 8 & 9 & 10 \\ \hline
        OPC & 0.00 & 2.50 & 1.25 & 0.00 & 0.00 & 0.00 & 1.25 & 0.83 & 1.67 & 0.42 & 0.42 \\ \hline
        Lime & 0.00 & 0.00 & 0.00 & 0.00 & 1.25 & 2.50 & 1.25 & 0.83 & 0.42 & 0.42 & 1.67 \\ \hline
        GGBFS & 0.00 & 0.00 & 1.25 & 2.50 & 1.25 & 0.00 & 0.00 & 0.83 & 0.42 & 1.67 & 0.42 \\ \hline
         \rowcolor{Goldenrod}
        Soil type B & ~ & ~ & ~ & ~ & ~ & ~ & ~ & ~ & ~ & ~ & ~ \\ \hline
        Recipe & 0 & 1 & 2 & 3 & 4 & 5 & 6 & 7 & 8 & 9 & 10 \\ \hline
        OPC & 0.00 & 2.50 & 1.25 & 0.00 & 0.00 & 0.00 & 1.25 & 0.83 & 1.67 & 0.42 & 0.42 \\ \hline
        Lime & 0.00 & 0.00 & 0.00 & 0.00 & 1.25 & 2.50 & 1.25 & 0.83 & 0.42 & 0.42 & 1.67 \\ \hline
        GGBFS & 0.00 & 0.00 & 1.25 & 2.50 & 1.25 & 0.00 & 0.00 & 0.83 & 0.42 & 1.67 & 0.42 \\ \hline
         \rowcolor{Goldenrod}
        Soil type C & ~ & ~ & ~ & ~ & ~ & ~ & ~ & ~ & ~ & ~ & ~ \\ \hline
        Recipe & 0 & 1 & 2 & 3 & 4 & 5 & 6 & 7 & 8 & 9 & 10 \\ \hline
        OPC & 0.00 & 2.50 & 1.25 & 0.00 & 0.00 & 0.00 & 1.25 & – & – & – & – \\ \hline
        Lime & 0.00 & 0.00 & 0.00 & 0.00 & 1.25 & 2.50 & 1.25 & – & – & – & – \\ \hline
        GGBFS & 0.00 & 0.00 & 1.25 & 2.50 & 1.25 & 0.00 & 0.00 & – & – & – & – \\ \hline
    \end{tabular}
    \label{tab:1}
\end{table}

Notation for Table \ref{tab:1}: Grain size distribution, particle density, natural water content and X-ray diffraction tests have been performed for all the soil types recipe 0. Optimum moisture content was tested for soil type A recipe 0 and type B recipes 0 and 1. Liquid limit was tested for soil types A and B, recipe 0. The UCS tests were performed on days 7 (soil types A and B, recipes 0 to 6, twice in each case), 28 (A: recipes 0 to 6, triple; B: all the recipes, twice; C: recipes 1 to 6, twice) and 90 (twice for all cases for A: all recipes; B: recipes 0 to 6; C: 0 to 7).

We estimated the mineral composition and structure of the solid phases of soil types A, B and C using the difference in X-ray diffraction (XRD), to understand its response on addition of binders during stabilization process. The principle consists in the adsorption of binders during reactions with soil, and transport of the materials of stabilizing agents through pores of soil, which performs at interfaces of soil with various physico-chemical properties that are different for types A, B and C in soil components. Soil types A and B contained clay minerals in a mixed layer of smectite (montmorillonite) \ch{(Na,Ca)0.33(Al,Mg)2(Si4O10)(OH)2*nH2O}, illite \ch{(K,H3O)(Al,Mg,Fe)2(Si,Al)4O10[(OH)2,(H2O)]}, kaolinite \ch{Al2(OH)4Si2O5}. Soil type C contains chlorite \ch{ClO-2}, illite \ch{K,H3O)(Al,Mg,Fe)2}\\\ch{(Si,Al)4O10[(OH)2,(H2O)]} and kaolinite \ch{Al2(OH)4Si2O5}. 


\subsection{Binders}
Binders can be rated according to their effects on strength development in soil (Table \ref{tab:2}). For example, the suggested binders for illite and kaolinite include sand, cement and lime. This explained the choice of lime and cement as binders for stabilization of soil type A and B. The chlorite component in soil type C well responds to treatment by cement, which is why the OPC was selected for soil type C. For these reasons, the following commercially available binders were used for soil treatment: ordinary Portland cement (OPC), ground granulated blast furnace slag (GGBFS) and a slaked lime. 

\begin{table}[!ht]
\caption{Suitability of binders in road stabilisation to soil types (Unified Soil Classification System, USCS)}
\small
    \centering
    \begin{tabular}{|l|l|l|l|l|l|l|}
    \hline
     	\rowcolor{Goldenrod}
        Stabilizing agent*** & G & GW & GM & S & CS & CH \\ \hline
        A & ++ & ++ & ++ & + & + & – \\ \hline
        B & ++ & ++ & ++ & ++ & ++ & + \\ \hline
        C & ++ & ++ & ++ & ++ & ++ & – \\ \hline
        D & + & + & ++ & – & + & ++ \\ \hline
        C (CaO + GP cement) & – & – & + & – & + & ++ \\ \hline
        C (CaO + GGBFS) & – & ++ & ++ & – & + & + \\ \hline
        E & ++ & ++ & + & + & – & – \\ \hline
        F & ++ & ++ & + & + & – & – \\ \hline
        G & + & ++ & ++ & – & ++ & + \\ \hline
    \end{tabular}
    \label{tab:2}
\end{table}

Notations for Table \ref{tab:2}: Definitions of symbols, according to USCS: G: gravel (crushed rock); GW: well-graded gravel; GM: Silty/clayey gravel; S: sand; CS: Sandy/silty clays; CH: fat heavy clay.
** Notations for suitability: ++ perfectly suitable; + satisfactory; – not suitable.
***Symbols for binders: A – GP cement AS3972; B – GB cement AS3972; C – Cementitious blend of pozzolanic binders (fly ash + GP cement + GGBFS + CaO in triple and quaternary mixtures); D – slaked lime or calcium hydroxide (CaO); E – Asphalt / bitumen, standard AS2008; F – GP cement / asphalt; G – Insoluble polymers.

The blends of standardized binders were used in blending operation, which was performed in a mixing with soil, according to standards of SGI for soil composition and recommended binders, and controlled through records used for ternary plots for UCS estimation. The we selected different types of binders for each type of soil, because these binders react differently with soil particles. The calculated difference in reaction products from 6 different binder recipes, which are as follows (in kg): 1) OPC (100); 2) OPC (50)+GGBFS (50); 3) GGBFS (100); 4) GGBFS (50) + lime (50); 5) lime (100); 6) lime (50) + OPC (50), Fig. \ref{fig:01}. The soil was treated with these binders or their blends. The total amount of binder used was set to 2.5\% of the soils' dry density, because it enables to obtain an expected maximum UCS of 5 MPa after a curing period of 3 months. A UCS of 5 MPa is a high value for soil stabilization, but was obtained using 2.5\% binder.

\begin{figure}[H]
      	\centering
         		\includegraphics[width=8.0cm]{Fig_01.jpg}
     \caption{The effects from OPC and lime on soil particles (10 g binder per 100 g dry soil)} \label{fig:01}
\end{figure}

All three soil types A, B and C contained material varying from clay to coarse sand particles and therefore could be treated with lime and cement. There are differences in the effects from various binders (Table \ref{tab:3}). For instance, OPC quickly starts to react with soil particles just within 2 h after mixing with soil. Most of the reaction is finished in 2 months, Fig. \ref{fig:01}. In contrast, the effects from lime are more consequent and similar to a straight line. 

\begin{table}[H]
\caption{Reaction on stabilisation from main soil compounds and suggested binders}
    \centering
    \begin{tabular*}{\textwidth}{c @{\extracolsep{\fill}} |l|l|l|}
    \hline
    \rowcolor{Goldenrod}
        Soil component & Chemical formula & Binders \\ \hline
        Organic matter & – & Mechanical \\ \hline
        Sand & \ch{SiO2} & Clay loam; asphalt \\ \hline
        Allophane & \ch{Al2O3·(SiO2)1.3-2·(2.5-3)H2O} & Lime (CaO) \\ \hline
        Kaolinite & \ch{Al2(OH)4Si2O5} & Sand; cement; lime \\ \hline
        Illite & \ch{K,H3O(Al,Mg,Fe)2(Si,Al)4O10[(OH)2,(H2O)} & Cement; lime \\ \hline
        Montmorillonite & \ch{(Na,Ca)0.33(Al,Mg)2(Si4O10)(OH)2·nH2O} & Lime (CaO) \\ \hline
    \end{tabular*}
    \label{tab:3}
\end{table}

Argumentation for the use of these materials is as follows. Organic matter increases in soil compaction due to molecular bonds between the particles. Sand contributes to the increase of the mechanical stability, density, cohesion and gain in shear strength of soil. Allophane contributes to the increase of the pozzolanic strength, compaction and densification of soil. Kaolinite has effects on the mechanical stability of soil, contributes to the early strength and workability as well as later strength gain of soil. Illite adds early strength gain and increases the workability and early strength development in soil. Finally, the montmorillonite effects the workability and early strength gain of soil.

To correspond between the long-term soil performance and workflow effectiveness in the geotechnical works, various manufacturers are producing specific binders with trade names. However, often the exact content and the proportions might not be available, which requires testing binders for soil stabilization. For example, lime and OPC mixed with clay produce reaction products after 18 months. However, lime-clay has more productivity compared to cement-clay with regard to reaction products, which is dependent on the temperature in the stabilized soil. OPC works better for coarse material, while lime is better suitable for fine-grained soil. Again, contrary to lime, OPC does not require external materials to result in a binding blend. The coarse silt type of soil can be stabilized by lime and cement better than by the pure binder.

\subsection{Compaction}
The samples were treated by vibratory surface compaction using the equipment available at SGI to compact the granular material of soil. The 4500 g of each soil specimens were compacted in three layers in a plastic tube with an outer steel tube. The surface of each compacted layer was scratched before the adding of the next layer. All the samples were compacted one hour after mixing with binders which excluded the evaluation of curing period from various binders. After compaction, the specimens were cut to a height of ca. 2.06 m, sealed with paraffin and stored within the plastic tube. 

The specimens did not lost weight during the curing period. Each layer was compacted with a 4-5 kg rammer with 25 blows, which resulted in decrease of 30 cm. To sift out large lumps, the soil passed through a 19 mm sieve and lumps larger than 19 mm were removed. The suitable specimens were then moulded again and placed into a sealed container, one by one. The containers were stored in a climate room at a constant temperature of 20°C and a constant relative humidity (RH) of 85\%. The MCV of the mixture, was determined in one hour of soil compaction, while the UCS tests were taken on days 7, 28 and 90th, respectively. 

\subsection{Moisture Condition Value (MCV)}
The MCV of the soil-binder mixture was performed because the compaction is a key procedure of soil improvement with regard to strength and durability. The advantage of MCV consists in its simplicity and availability to measure the lowest compaction energy needed to achieve the maximum soil density at a specific water content for assessing the fitness of soil for geotechnical road works. The MCV was determined in accordance to the standards EN 13286-46 and the UK Standard 1377 \cite{SIS2003,BIS}, adapted for current case. This MCV is based on the repeated compaction of specimen with rammer until the limited compaction is reached. 

The technical procedure is as follows. A free-falling 7 kg and 97 mm-diameter rammer was attached to an automatic release device. A soil sample of 1.5 kg was used with a drop height of 2.50 m. A lightweight disk was placed on top of the soil to prevent extrusion of soil and smearing of the rammer. The upper limit for MCV was set up at 14, to ensure the optimum moisture content for stabilised soil. The dry densities of soil specimens, stabilised by 6 types of binders, described in previous section were compacted by MCA. 

The results are compared with the dry density of the vibratory-compacted samples for soil types A and B, Fig. \ref{fig:02}. The dry density is higher for type A compacted by MCA, compared to the vibratory compacted, Fig. \ref{fig:02}. In contrast, soil type B demonstrated the same or lower dry density values for the MCA-compaction, compared to the vibratory compaction. The MCV for soil A is slightly higher than that for type B, due to the techniques of the vibratory compaction which does not produce dry density similar to that in MCA compaction for soil A. 

\begin{figure}[H]
      	\centering
         		\includegraphics[width=14.0cm]{Fig_02.jpg}
     \caption{Dry density for different compaction approaches and binder blends for soil types A (left) and B (right).} \label{fig:02}
\end{figure}

Nevertheless, the vibratory compaction results in identical dry density as the MCA. The MCV values are below 14 for all tested samples. 
\begin{figure}[H]
      	\centering
         		\includegraphics[width=7.0cm]{Fig_03.jpg}
     \caption{Natural water content for soil types: (a) A; (b) B; and (c) C, respectively.} \label{fig:03}
\end{figure}

For the case of MCV <14, the soil is workable and becomes homogeneous and dense with normal compaction equipment.  Each soil type was homogenized, however, small variations in the water content remained between each specimen due to the individual properties of soil, which lied within normal distribution.
Water content in the stabilized soil is an important variable for groundwater discharge and evaluation of soil chemistry. Thus, if the moisture content of a soil is optimum, it can be suitable for construction works. At the same time, due to the variations in the mineralogical content (clay, quartz, etc.), natural water content differs in soil types A, B and C within the same group due to the individual properties, and between the investigated soil types, Fig. \ref{fig:03}.  

\subsection{Simplex Lattice Design}
The mixture design was adopted from existing studies to evaluate the effects from individual binders and their blends using simplex lattice design with 3 binders: OPC, GGBFS and slaked lime. We used six components of recipes for normal simplex lattice design. The scheme presents three single blends of 100\% of a single ingredient, and three mixtures of two of the binders (Fig. \ref{fig:04}). 

\begin{figure}[H]
      	\centering
         		\includegraphics[width=7.0cm]{Fig_04.jpg}
     \caption{Simplex lattice design for binder selection (3, 2)} \label{fig:04}
\end{figure}

Three statistical regression models model the material, depending on the test design. These include linear, quadratic and special cubic models. 
(a) Linear model:
(b) Quadratic model:
(c) Special cubic model:

where $\beta$ is the regression coefficient; $\epsilon$ is the residual error; $x_i$ is the value of factors; $y$ is the dependent variable.

The dependent variable y contains effect $x_1<\ldots>x_3$ (linear model), interactions $x_1x_2$, $x_1x_3$, $x_2x_3$ (quadratic model) and $x_1x_2x_3$ (special cubic model). We included only the significant effects from the quadratic and special cubic regression models. The regression models were tested with the ANOVA to evaluate the significance. For models with not significant ANOVA, each effect was tested for significance and those with a p>0.05 were removed.  The null hypothesis $H_0$ was that the mean of MCV or UCS of the stabilized soil was equal for all binders. This alternative hypothesis $H_1$ is that at least one binder differs from the others. Testing the equality of binders leads to the equation (4): 

\begin{equation}
H_0 = \mu _1 = \mu _2 = <\ldots> = \mu _a
\end{equation}

\begin{equation}
H_1: \mu _i \neq \mu _j
\end{equation}
for at least one pair of $(i, j)$.

If $H_0$ is true, then all treatments have a common mean $\mu$ [52]. The p-level indicates the probability of rejecting the $H_0$ when it is actually true. The $H_0$ was statistically tested and rejected at a significance level (p=0.05). 

\section{Results and Discussion}
Results demonstrate a series of plotted ternary plots showing response surfaces for UCS and MCV (Figs. 5 to 8). The circles on the borders of these figures signify a mixture of single binder, binary blend or a complex mixture. The results of the MCV tests are presented for soil types A and B in Fig. \ref{fig:05}. Straight lines representing the contours signify no interaction between the binders as independent variables. In other words, no 2nd-order term is significant (equations 1 and 2 for linear and quadratic models). Response surfaces in Fig. \ref{fig:05} to \ref{fig:08} is show the effects from binders on soil stabilization. 

\begin{figure}[H]
      	\centering
         		\includegraphics[width=14.0cm]{Fig_05.jpg}
     \caption{MCV for various binders blended with soil: (a) type A; (b) type B} \label{fig:05}
\end{figure}

\subsection{MCV}
The response surface equation for binder blends tested for MCV on soil types A and B is presented in equations (6) and (7). Here, C, L and S stand for cement, lime and steel slag, and combinations LS and CLS – for their mixes, correspondently. The numbers represent the amount of binders taken for stabilization.

MCV=13.2C+11.4L+8.9S+3.7CS+2.2LS 		(6) 
MCV=12.4C+10.5L+8.6S+7.4CL+22.5CLS	 (7)

The total amount is adjusted for 100\%, that is, boundary conditions satisfy to the following assumption:  $0 \le C$, $S, L \le 1$ and $C+L+S = 1$. The sum of the independent variables is 1 for all cases. One can derive from equations (6) and (7), that cement has the greatest impact on soil compaction after one hour of treatment time. There are correlations between various binders for all the soil types with those having a positive mark indicating the increase in MCV.

Nevertheless, the contribution to the MCV can take up to 0.925 at maximum for soil A with a blend OPS/GGBFS 50/50\% (eq. 6). The contribution to the MCV reaches 1.85 at maximum for soil B for blend of binders OPC/lime 50/50\% (eq. 7).  For soil type B, the MCV relationships resulted in both positive and negative cases. Binder blends gave a significant contribution to the MCV. For soil type A, there is an adjusted R2 value of 0.999, i.e., a very low variability. GGBFS gives the smallest effect on soil compaction for both soil types, Fig. \ref{fig:05}

\subsection{UCS}
The UCS tests were performed for soil types A and B on days 7, 28 and 90th, while soil type C was tested on days 28 and 90th. 

\subsubsection{Measurements of UCS on 7th day}
Because soil types varied in clay content, we assumed that lime would perform differently in soil types A–C, which is proved by Fig. 6 to 8. 
UCS=+2930C+850L+310S–1640CL +2500CS+2150LS (8) 
UCS=+1500C+640L+240S +1650LS (9)
Equations (8) and (9) represent measured 7-day UCS (kPa) for soil types A and B.

\begin{figure}[H]
      	\centering
         	\includegraphics[width=14.0cm]{Fig_06.jpg}
     \caption{UCS strength for soil-binder mixtures cured for 7 days: (a) A; (b) B} \label{fig:06}
\end{figure}

Fig. \ref{fig:06} shows the UCS for soil A and B treated during 7 days. Soil type A demonstrated a higher strength compared to the types B and C. Besides, it also has a significantly larger UCS interval compared to type B. Indeed, soil type A has the UCS varying within 2-6 MPa, while UCS of type B varies from 0.8-2.7 MPa. For soil type A, the 7-day UCS shows significant full quadratic regression model, i.e., all effects from all binders were considerable. We also noted that OPC has the largest impact on the UCS of soil, while the GGBFS has the lowest impact on UCS for both soil types, Fig. \ref{fig:06}. After a 7-day of curing period, the OPC-lime blend has a negative effect on UCS, i.e., that this part of the blend reduces the UCS. For soil type B, the UCS on 7th day of curing results in linear model and has quadratic significant term lime-GGBFS. With a 50:50\% blend of lime-GGBFS, the quadratic term has more impacts on the UCS compared to the effects from single slaked lime. 

\subsubsection{Measurements of UCS on 28th day}
The UCS of stabilized soil of types A, B and C, on 28th days of curing period are presented in the subplots of Fig. \ref{fig:07}. The UCS of soil in kPa for soil type A, B and C on 28th day of curing is summarized in equations (10), (11) and (12). Here C, L and S indicate OPC, lime and GGBFS, and combinations CS and LS – for their mixes, respectively.
UCS=+3800C+1030L +1210S+4270CS+6030LS (10) 
UCS=+2080C+810L+410S+4440LS	(11)
UCS=+1090C+180L+430S+2360CS	(12)

The OPC-lime blend for soil type A was not significant and removed from the regression equation (eq. 10).  On day 28th of curing OPC remains the dominating element on the UCS of soil. However, the GGBFS is larger compared to lime component, eq. (8). For soil type B, samples treated for 28 days (eq. 11) received significant effects the same as those cured for 7 days (eq. 9). Nevertheless, the regression coefficient for the slaked lime/steel slag blend increased more distinct compared to the other blends. This suggests that lime is well suited to soil type B as a binder. The check points were added to validate the results in the UCS test with 28 days of curing period, which show that the adjusted R2 value is satisfactory. The effects from lime is not significant for soil type C. The contribution from the lime effect in kPa (eq. 12) is smaller, as to soil types A and B and can be considered as minor. Such minor effect from the slaked lime is owing to the low clay content in soil type C. Another factors can be explained by short curing time, while the effects from lime are time-dependent, see Fig. \ref{fig:01}. Certainly, slaked lime requires a longer curing time to result in reaction similar to that from cement and steel slag.  

\begin{figure}[H]
      	\centering
         	\includegraphics[width=7.0cm]{Fig_07.jpg}
     \caption{UCS strength for soil-binder mixtures cured for 28 days: (a) A; (b) B; (c) C.} \label{fig:07}
\end{figure}

\subsubsection{Measurements of UCS on 90th day}
The UCS of soil samples cured for 90 days are presented in Fig. \ref{fig:8}. For soil type A, the regression model is shown in eq. 13 with the least effective binder component as OPC-GGBFS. This can be compared with eq. 8 and 10 showing significant effects from OPC-GGBFS blend. The notable correlation after 90 days curing period is the lime-GGBFS blend. The effect of the lime-GGBFS variable with a 50/50\% percentage is greater compared to single binder of slaked lime.

\begin{figure}[H]
      	\centering
         		\includegraphics[width=7.0cm]{Fig_08.jpg}
     \caption{UCS for soil types: (a) A; (b) B; (c) C cured for 90 days with various binders.} \label{fig:08}
\end{figure}

The effects from blends of OPC/lime and lime/GGBFS are shown for soil type B in eq. 14. The maximum UCS lie in a range between OPC/lime and lime/GGBFS blends (Fig. \ref{fig:8}). As for soil type C, the significant effects have linear character with GGBFS giving the highest UCS after a curing period of 90 days (Fig. \ref{fig:8} and eq. 14). 
	UCS=+5880C+2270L+ 2960S+4600LS (13)
	UCS=+2390C+2140L+810S+1390CL+4290LS (14) 
	UCS=+1340C+371L+1910S (15)
 The clay content in soil type C is lower than in soil types A and B (6\%), which suggests that slaked lime has a lesser effect on the UCS of soil type C compared to types A and B. The minimum 7-day cube UCS is specified for cement-bound material 1 (CBM 1) and cement bound material 3 (CBM 3), Table \ref{tab:4}. Here, none of the binders reach the limits for CBM 1-3 with 2.5\% of binder content. 
 
 \begin{table}[H]
 \caption{Requirements for compressive strength (MPa) for soil stabilised by cement after 7 days of curing}
    \centering
    \begin{tabular}{|l|l|l|l|}
    \hline
       \rowcolor{Goldenrod}
        Shape form & Category 1 & Category 2 & Category 3 \\ \hline
        Cube specimens & 4.5 & 7.5 & 10.0 \\ \hline
        Cylindrical specimens & 3.6 & 6.0 & 8.0 \\ \hline
    \end{tabular}
    \label{tab:4}
\end{table}
 
 The OPC-GGBFS and GGBFS result in soil UCS <2 MPa after a 90-day curing period. This is because the quantity of calcium hydroxide in these blends is not sufficient to form a uniform mixture which would produce pozzolanic material. The GGBFS-lime blend shows the most notable UCS increase in the period between 7 and 28 days. Table 5 shows that a mixture of OPC and lime gives higher UCS compared to pure OPC or lime. The clay content in the sandy clay was 17\%, versus clay content in control clay (55\%). It means that the results of soil stabilization additionally dependent on clay content.

%%%%%%%%%%%%%%%%%%%%%%%%%%%%%%%%%%%%%%%%%%
\section{Conclusion}

The strength properties of soil types A and B (sandy silty tills), and C (clay till) were evaluated using selected blended binders taken as pure agents and in mixes. We tested ordinary Portland cement (OPC), ground granulated blast furnace slag (GGBFS) and a slaked lime and their blended mixtures to evaluate the Uniaxial Compressive Strength (UCS) and Moisture Condition Value (MCV) in treated soil, examined the mineral structural and content properties of the soil using X-ray diffraction (XRD) and preliminary sieve and assessed the water content. The tests were intended to improving the properties of soil prior to road construction. The obtained results led to the following conclusions:

\begin{enumerate}
   \item The strength increase is comparable for most of the tested binders with notable effects from cement which works rapidly and results in quick soil stabilization.
   \item The UCS increase in soil stabilized by OPC-lime blend is similar to that of GGBFS. However, a binder containing GGBFS gives lower values.  
   \item The measurements on day \nth{7} of curing show that the mixture of GGBFS and slaked lime performs similar to the pure slaked lime. However,  after a curing period of 90 days there is a notable benefit for soil strength from the GGBFS-lime blend. 
   \item The specimens stabilized with GGBFS show the lowest UCS after \nth{7} day of curing time. However, on \nth{28} day of curing they perform similar to those stabilized by lime.
   \item Low increase in UCS for the lime-stabilized specimens is explained by low clay content in soil type A. slaked lime performs better as a binder, compared to others, and has a stronger response in soil type B compared to type A.
   \item The GGBFS-lime blend demonstrated benefit on UCS of soil and performs approximately in similar way as OPC for soil strength after \nth{90} day of curing.	
   \item A lower limit for clay content in soil should be ca. 10\% for lime binder to be effective.
   \item We did not find a significant variations in effects from blending the three binders with regard to reduce the soil compaction, as reflected in the MCV values
 \end{enumerate}

	This article demonstrated the effects from single binders (cement/lime/steel slag) and their blended mixes on various types of clayey soils collected in southern Sweden. Although traditional methods of soil stabilization widely apply OPC or slaked lime as stabilising agents, the increase in using blended binders is notably recently. This is caused by the improved properties of such mixtures. As a result, blending OPC with GGBFS or OPC with pozzolanic materials, such as fly ash, has gained attention in civil engineering works. The vibratory compaction method that we applied in this experiment results in similar dry density values of soil if compared to the specimens compacted by MCA. Therefore, it could be possible to achieve these values in-situ, as the maximum MCV was lower that 14.
	We noted positive effects from the blends GGBFS/lime on UCS gain in soils type A and B. Moreover, the GGBFS/lime mixed binder performs better compared to pure GGBFS or lime. In soil type B, the blend of GGBFS and slaked lime even outperforms pure OPC, owing to the fact that blend of GGBFS/lime does not affect compatibility of soil as a OPC binder, inasmuch as all the specimens were produced with the identical compaction approach. Second, clay content in soil type B is enough to produce maximum reaction products, while in soil type A, it is too low to match the reaction products produced by OPC binder.
	The main aim of blending stabilizing agents is to increase the economic effectivity of workflow, following by environmental goals of utilization of waste material, and experimental producing of a new binder with better properties for technical reasons. Nonetheless, blending binders does not mechanically leads to the improved properties of stabilizing agent. Instead, the opposite reaction may be as the result depending on various factors, including soil grain size, or temperature among the others. Therefore, mixed binders were tested in this study, to ensure that they perform better than single binders for a particular soil from southern Sweden.

%%%%%%%%%%%%%%%%%%%%%%%%%%%%%%%%%%%%%%%%%%
\section{Conclusions}

Technical implementation of soil stabilization and compaction is very expensive. At the same time, such works are an absolute requirement prior to all construction projects, including construction of buildings, foundations, infrastructure, tunnels, roads, railways and highways, and associated transport infrastructure. Indeed, soil quality is one of the crucial factors in civil engineering which ensures safety and workability of constructed objects. Therefore, selecting advanced, robust, reliable and optimized methods for soil compaction and stabilization are essential for industrial purposes, as compaction and UCS are perfect indicators of soil quality, achieved through successfully performed stabilization.

Advanced statistical techniques need to be used to optimise the ratio of binders aimed at determining the effects of their blended mixtures. Here we presented a series of simplex tests used to visualize and plot the effects from stabilizing agents on clayey soil. This test is beneficial for testing one amount is tested per time for determining the exact amount of single binders to be added. For more complex cases, it is recommended to perform several ternary plots for estimation of various amounts of binders, as independent statistical variables. However, in this case, the total amount of runs requires much greater trials which results in a more complicated of assessment soil strength stabilized by various from binders. 

Correctly selected binders, adjusted to the specific type of soil, result in a better stabilised soil that performs well. Moreover, well-adjusted techniques can significantly improve time-consuming workflow time and reduce economic costs of project implementation. For example, in the demonstrated study using pure OPC for stabilization of soil types A and B resulted in the best effects on MCV. Soil stabilized by the OPC needs a higher compaction energy due to the much higher their stiffness. 
	
The limitations of the presented work include the scope of the research, which is limited to the values received in the laboratory. Therefore, we cannot directly compare the results of the tested soil with those received on soil tested in real field in-situ conditions. At the same time, real conditions may affect soil response due to the regional and local properties and environmental-climate setting. As a result, this may lead to the variations in physico-chemical parameters of soil and mixing conditions. As a recommendation for future similar works, testing different various types of blends (proportions of binders within the blends) and different amounts of binders (in absolute mass, kg) would be beneficial. The improvement of soil compaction in situ should ideally include investigation of responses from compacted soil and the soil underneath on compaction. 

%%%%%%%%%%%%%%%%%%%%%%%%%%%%%%%%%%%%%%%%%%
\vspace{6pt} 

%%%%%%%%%%%%%%%%%%%%%%%%%%%%%%%%%%%%%%%%%%
\authorcontributions{For research articles with several authors, a short paragraph specifying their individual contributions must be provided. Conceptualization, methodology, software, validation, visualization, supervision -- Per Lindh; formal analysis, investigation, resources, data curation, project administration, funding acquisition -- Polina Lemenkova; writing---original draft preparation, Polina Lemenkova ; writing---review and editing, Per Lindh and Polina Lemenkova. All authors have read and agreed to the published version of the manuscript.}

\funding{This research was funded by Electronics Editorial Office, MDPI, providing 100\% discount for APC for publication of this manuscript.}

\institutionalreview{Not applicable}

\informedconsent{Not applicable}

\dataavailability{Not applicable} 

\acknowledgments{We thank the two anonymous reviewers for comments and reading of this manuscript. The study was supported by the NUTEK – Swedish Agency for Economic \& Regional Growth, Peab AB construction and civil engineering enterprise and Merox. We thank Cementa (HeidelbergCement Group Northern Europe) for provided cement binder, and Nordkalk for limestone binder. The authors acknowledge the help from colleagues from the Lund University, Chalmers Institute of Technology and the Swedish Geotechnical Institute with laboratory assistance and advices regarding the experimental tests. }

\conflictsofinterest{The authors declare no conflict of interest} 

\abbreviations{Abbreviations}{
The following abbreviations are used in this manuscript:\\

\noindent 
\begin{tabular}{@{}ll}
ANOVA & Analysis of Variance \\
GGBS & Ground Granulated Blast Furnace Slag \\
MCV & Moisture Condition Value \\
OMC & Optimum Moisture Content \\
OPC & ordinary Portland cement \\
SGI & Swedish Geotechnical Institute \\
UCS & Uniaxial Compressive Strength \\ 
WN & natural water content \\
XRD & X-ray diffraction
\end{tabular}
}


%%%%%%%%%%%%%%%%%%%%%%%%%%%%%%%%%%%%%%%%%%
\begin{adjustwidth}{-\extralength}{0cm}
%\printendnotes[custom] % Un-comment to print a list of endnotes

\reftitle{References}

% Please provide either the correct journal abbreviation (e.g. according to the “List of Title Word Abbreviations” http://www.issn.org/services/online-services/access-to-the-ltwa/) or the full name of the journal.
% Citations and References in Supplementary files are permitted provided that they also appear in the reference list here. 

\bibliography{SimplexLattice.bib}
\nocite{*}

%%%%%%%%%%%%%%%%%%%%%%%%%%%%%%%%%%%%%%%%%%
\end{adjustwidth}
\end{document}

